\setcounter{ex}{0} %Ba Đồn-Quảng Bình
\subsection{Trắc nghiệm (2 điểm)}
\Opensolutionfile{ans}[ans/ans7]

%%Câu 1
\begin{ex}
Phương trình nào sau đây là phương trình bậc nhất hai ẩn?
\choice
{$2\sqrt{x}+2=0$}
{$2y-1=2(y+1)$}
{\True $2x+\dfrac{y}{2}=1$}
{$x+y^2=0$}
\end{ex}

%%Câu 2
\begin{ex}
Cho biết số nghiệm của hệ phương trình $\heva{&y=2x-1\\&y=-x+2}$ là:
\choice
{Vô nghiệm}
{Có hai nghiệm}
{Có vô số nghiệm}
{\True Có một nghiệm duy nhất}
\end{ex}

%%Câu 3
\begin{ex}
 Điều kiện để biểu thức $\dfrac{x-2}{2x-4}\cdot \dfrac{\sqrt{x}-1}{\sqrt{x}+2}$ xác định
\choice
{\True $x\ge 0, x\ne 2$}
{$x\ge 0$}
{$x\ge 2$}
{$0\le x\le 2$}
\end{ex}

%%Câu 4
\begin{ex}
 Với mọi số thực $a,b$ thỏa mãn $a > b$, khẳng định nào sau đây là đúng?
\choice
{$\sqrt{a}>\sqrt{b}$}
{\True $a>\dfrac{a+b}{2}$}
{$3a>-3b$}
{$a<2b$}
\end{ex}

%%Câu 5
\begin{ex}
Rút gọn biểu thức $\dfrac{\sqrt{x^2-2x+1}}{\sqrt{x}-1}$ với $0\le x<1$, ta được:
\choice
{$1+\sqrt{x}$}
{$\sqrt{x}-1$}
{$1-\sqrt{x}$}
{\True $-1-\sqrt{x}$}
\end{ex}

%%Câu 6
\begin{ex}
 Cho tam giác $ABC$ là tam giác đều nội tiếp đường tròn $(O)$ có bán kính $R = 2$. Kẻ đường kính $BD$ của đường tròn $(O)$. Khi đó độ dài cạnh $CD$ bằng
\choice
{4}
{$2\sqrt{3}$}
{\True 2}
{$2\sqrt{2}$}
\end{ex}

%%Câu 7
\begin{ex}
Cho tam giác $ABC$ vuông tại $A$ có cạnh huyền $BC=a$, cạnh góc vuông $AC=b$, $AB=c$. Ta có:
\choice
{$b=c\cdot\sin\widehat{B}$}
{\True $b=a\cdot\cos\widehat{C}$}
{$c=b\cdot\sin\widehat{B}$}
{$c=b\cdot\cos\widehat{B}$}
\end{ex}

%%Câu 8
\begin{ex}
Số tâm đối xứng và trục đối xứng của đường tròn $(O;R)$ là:
\choice
{1 tâm đối xứng và 1 trục đối xứng}
{Vô số tâm đối xứng và 1 trục đối xứng}
{\True 1 tâm đối xứng và vô số trục đối xứng}
{Không có tâm đối xứng và trục đối xứng}
\end{ex}

%Câu 9
\begin{ex}
Cho tam giác $ABC$ nội tiếp đường tròn $(O)$. Biết rằng $\widehat{BAC}=15^\circ$◦. Tính $\widehat{BOC}$
\choice
{\True $30^\circ$}
{$45^\circ$}
{$60^\circ$}
{$90^\circ$}
\end{ex}

%Câu 10
\begin{ex}
Từ một điểm $M$ nằm ngoài đường tròn $(O;R)$ vẽ các tiếp tuyến $MA,MB$ với đường tròn ($A,B$ là các tiếp điểm). Biết rằng diện tích tứ giác $OAMB$ là $\sqrt{3}R^2$. Khi đó $AB$ bằng
\choice
{$2R$}
{\True $\sqrt{3}R$}
{$\sqrt{2}R$}
{$R$}
\end{ex}

\Closesolutionfile{ans}
\Opensolutionfile{ansbook}[ans/ansbook7]
\setcounter{bt}{0}
\subsection{Tự luận (8 điểm)}
%Bài 1
\begin{bt}[1,5 điểm]
Giải các phương trình và hệ phương trình sau:
\begin{listEX}[2]
\item $(x-1)(3x-6)=0$
\item $\heva{&2x+2y=9\\&2x-3y=4}$
\end{listEX}
\loigiai{ 
a) $x=1, x=2$ \hspace{2cm } b) $\left(\dfrac{7}{2};1\right)$
}
\end{bt}
%Bài 2
\begin{bt}[1,5 điểm]
Cho biểu thức 
$$P=\left(\dfrac{1}{\sqrt{a}-1}-\dfrac{1}{\sqrt{a}}\right):\left(\dfrac{\sqrt{a}+1}{\sqrt{a}-2}-\dfrac{\sqrt{a}+2}{\sqrt{a}-1}\right)\ \text{(với } a>0,\, a\ne1,\, a\ne4)$$
\begin{enumerate}
\item Rút gọn $P$.
\item So sánh $P$ với $\dfrac{1}{3}$.
\end{enumerate}
\loigiai{
a) $P=\dfrac{\sqrt{a}-2}{3\sqrt{a}}$ \hspace{2cm} b) $P<\dfrac{1}{3}$
}
\end{bt}
%Bài 3
\begin{bt}[1,0 điểm]
 Một chiếc ô tô dự định di chuyển từ địa điểm A đến địa điểm B để du lịch. Theo tính toán ban đầu, xe định đến B lúc 10h sáng và nếu mỗi giờ xe chạy chậm hơn dự kiến 10 km thì sẽ đến nơi lúc 3h chiều. Tuy nhiên, do gặp vấn đề về sự cố nên xe đã xuất phát trễ hơn dự kiến 3h, do đó để kịp đến nơi như dự định thì vận tốc của xe cần tăng thêm 10 km/h. Hỏi quãng đường từ A đến B là bao nhiêu kilomet và xe đã xuất phát lúc mấy giờ?
\loigiai{
Gọi $x$ (km/h) và $y$ (h) lần lượt là vận tốc dự định và thời gian dự định.

Quãng đường AB là 600 km, xe xuất phát lúc 7 giờ tối hôm trước.
}
\end{bt}
%Bài 4
\begin{bt}[3,0 điểm]
Cho đường tròn $(O)$ đường kính $AB$. Lấy $M$ thuộc $(O)$ sao cho $MA<MB$. Vẽ dây $MN$ vuông góc với $AB$ tại $H$. Đường thẳng $AN$ cắt $BM$ tại $C$. Đường thẳng qua $C$ vuông góc với $AB$ tại $K$ và cắt $BN$ tại $D$.
\begin{enumerate}
\item Chứng minh $A, M, C, K$ cùng thuộc một đường tròn.
\item Chứng minh $BK$ là tia phân giác của $\widehat{MBN}$.
\item Chứng minh tam giác $KMC$ cân và $KM$ là tiếp tuyến của đường tròn $(O)$.
\end{enumerate}
\loigiai{
\immini
{\begin{enumerate}
    \item $A,M,C,K$ thuộc đường tròn tâm $I$ bán kính $IM$.
    \item $\triangle MHB=\triangle NHB$, suy ra $\widehat{MBH}=\widehat{NBH}$.
    \item $\triangle KMH=\triangle KNH$, suy ra $KM=NK$

    $\widehat{KMO}=180^\circ-\widehat{KMC}-\widehat{OMB}=180-\widehat{KCM}-\widehat{OBM}=\widehat{CKB}=90^\circ$
\end{enumerate}}
{\begin{tikzpicture}[line cap=round,line join=round,font=\footnotesize]
		\def\r{2}
		\path
		(-\r,0) coordinate (A)
		(\r,0) coordinate (B)
		(0,0) coordinate (O)
		($(O)!1!-50:(A)$) coordinate (M)
		($(A)!(M)!(B)$) coordinate (H)
		($(M)!2!(H)$) coordinate (N)
		($(B)!2!(M)$) coordinate (Bt)
		($(N)!2.5!(A)$) coordinate (Nt)
		(intersection of N--Nt and B--Bt) coordinate (C)
		($(A)!(C)!(B)$) coordinate (K)
		($(C)!2!(K)$) coordinate (Ct)
		($(B)!1.5!(N)$) coordinate (Nt)
		(intersection of C--Ct and B--Nt) coordinate (D)
		($(C)!0.5!(A)$) coordinate (I)
		;
		\draw 
		(O) circle (\r)
		(M)--(N) (K)--(A)--(B)--(C)--(D)--(B) (C)--(N) (M)--(A)--(B)--cycle (M)--(I)--(K) (M)--(O)--(N)--(K)--(M)
		;
		\foreach \x/\g in{A/-135,B/-45,O/-90,M/90,H/-45,N/-90,C/90,K/180,D/-90,I/70}
		\draw[fill=white] (\x) circle (1.5pt) +(\g:.35) node{$\x$};
	\end{tikzpicture}}
}
\end{bt}
%Câu 5
\begin{bt}[1,0 điểm]
Một người gửi tiền vào ngân hàng với lãi suất $0{,}5\%$/tháng. Biết rằng, nếu không rút tiền ra khỏi ngân hàng thì cứ sau mỗi tháng, số tiền lãi sẽ được nhập vào vốn ban đầu. Người đó phải gửi số tiền ban đầu ít nhất là bao nhiêu triệu đồng để số tiền lãi sau tháng thứ hai không ít hơn $500\,000$ đồng (tính kết quả theo đơn vị triệu đồng và làm tròn kết quả đến hàng đơn vị)?
\loigiai{
$x$ (triệu đồng) là số tiền ban đầu người đó gửi vào ngân hàng

$0,5\%\left(1+0,5\%\right)\ge 0,5$

Người đó cần gửi số tiền ban đầu ít nhất là 100 triệu đồng. 
}
\end{bt}

\Closesolutionfile{ansbook}